%---
% slug: kpsa-class-sensitivity
% title: How ImageNet classes divide a vision transformer's parameter sensitivity
% category: update
% eyebrow: INTERPRETABILITY
% summary: We measure which MLP parameter groups the training loss of a supervised and a self-supervised ViT-B/16 is sensitive to on 50,000 ImageNet images. Sensitivity becomes class-specific in the second half of both networks, it is carried increasingly by WordNet supergroups with depth, and its class geometry follows the WordNet hierarchy most closely in the middle blocks.
% author: Davide Wiest
% draft: false
% listed: true
% date: 2026-11-17
% rank: 1
% cover_image: /work/covers/kpsa-class-sensitivity.svg
%---
\documentclass[11pt]{article}
\usepackage{publication-web}
\begin{document}

\publicationtitle{How ImageNet classes divide a vision transformer's parameter sensitivity}
\publicationauthors{Davide Wiest}
\publicationaffiliations{Axym Labs}

\begin{publicationhead}
\makepublicationtitle
\end{publicationhead}

A trained network does not use all of its parameters for every input. \href{https://axym.org/work/kpsa/}{Kernelized Parameter Sensitivity Atlases} (KPSA) build on this: they store where a model's sensitivity lies for a set of source inputs and retrieve it for new ones \cite{kpsa}. Retrieval works when inputs that are alike in a query space also rely on similar parameters. In this update we ask how this reliance is organized by the categories of the data, using ImageNet's 1,000 classes and the WordNet hierarchy above them.

We compare two ViT-B/16 models with identical architecture: one trained with labels on ImageNet-1k and one trained without labels by DINO self-distillation \cite{steiner2022,caron2021}. Across both, sensitivity is shared by all images in the first blocks and becomes class-specific in the second half of the network. With depth, a growing share of the parameter groups' class selectivity is selectivity for broad WordNet categories, and the class geometry of sensitivity follows the WordNet hierarchy most closely in blocks 6 and 7. A single image relies on an effective 6\% of the parameter groups in the supervised model and 13\% in the self-supervised one. The class-to-category structure connects to the \href{https://axym.org/work/hierarchical-data-hypothesis/}{hierarchical data hypothesis}, which states that broad categories capture what their subcategories share \cite{hdh}. You can explore input-conditioned sensitivity for single images and words in the \href{https://axym.org/work/kpsa-demo/}{interactive KPSA demo}.

\section{Setup}\label{setup}

For an input \(x\) with training loss \(\ell(x)\), the sensitivity of parameter group \(m\) is its squared gradient norm \(\lVert\nabla_{\theta_m}\ell(x)\rVert^2\), the diagonal block of the empirical neural tangent kernel. Normalized over all groups, it gives the profile \(s(x)\) that KPSA stores, a distribution over parameter groups \cite{kpsa}. The groups are the 36,864 MLP features of the 12 transformer blocks: each couples one row of the first MLP projection, its bias, and the matching column of the second projection. A group's gradient is a sum over tokens, so its energy splits exactly into per-token terms, which we place on the 14 by 14 image grid.

Each model is profiled under the loss it was trained on. For the supervised AugReg model this is the negative log-likelihood of the labelled class; it classifies 76.1\% of the images correctly. For DINO it is the self-distillation cross-entropy between the frozen teacher and the student over two global and ten local crops, drawn with a fixed seed per image, with the student's gradient. The released checkpoint does not contain the teacher's centering vector, so we use the mean teacher output over 20,000 crops of training images. Every analysis uses 50 validation images per class, 50,000 images per model.

Supergroups come from the WordNet hypernym tree: we split it from the root until every group holds at most 250 or at most 50 classes, which gives 35 and 155 supergroups. We compare profiles with the Bhattacharyya coefficient \(\sum_m \sqrt{s_m(x)\,s_m(x')}\), which is 1 for identical profiles and 0 for disjoint ones.

\section{How stable are sensitivities across instances of a class?}\label{stability}

In blocks 1 to 4, any two images have similar profiles, between 0.79 and 0.90, and pairs from one class are at most 0.03 more similar than pairs from different supergroups. From block 5 on, this gap widens. In the supervised model, two images of the same class reach 0.57 in block 11, compared with 0.14 for different classes of one supergroup and 0.08 for different supergroups. DINO separates the three kinds of pairs less sharply, with 0.33, 0.21 and 0.15 in block 11. Supergroup membership thus raises similarity in both models, even though DINO never saw a label.

\publicationfigure[\linewidth]
  {figures/kpsa-class-sensitivity-fig1-stability.png}
  {/work/kpsa-class-sensitivity/fig1-stability.svg}
  {Two line plots of profile similarity against transformer block for pairs of images of the same class, of the same supergroup, and of different supergroups}
  {Mean Bhattacharyya coefficient between the block-restricted profiles of two images of the same class, of different classes in one of 155 supergroups, and of different supergroups. The band spans the 25th to 75th percentile of the same-class value over the 1,000 classes.}

\section{How selective are parameter groups for classes and their supergroups?}\label{selectivity}

We measure a group's class selectivity as the information, in bits, that the distribution of its sensitivity over the classes carries relative to the class frequencies, weighted by the group's share of total sensitivity and with the value under shuffled image labels subtracted. Selectivity stays at or below 0.11 bits in blocks 1 to 4 of both models and rises with depth to 3.21 bits in block 11 of the supervised model and 1.80 bits in block 12 of DINO.

The same quantity computed over supergroups instead of classes, divided by the class value on the same classes, gives the share of a group's class selectivity that only distinguishes supergroups. The 35 coarse supergroups account for 37\% of the class selectivity in block 1 of the supervised model and for 58\% in block 10; in DINO the share rises from 35\% to 68\% in block 8. For the 155 supergroups, the share reaches 84\% and 92\%. Deeper parameter groups are therefore more selective, and most of their selectivity separates broad categories rather than single classes.

\publicationfigure[\linewidth]
  {figures/kpsa-class-sensitivity-fig2-selectivity.png}
  {/work/kpsa-class-sensitivity/fig2-selectivity.svg}
  {Three line plots: class information against block for both models, and the share of class information carried by supergroups against block for each model}
  {Left: mass-weighted information of parameter-group sensitivity about the class, above a shuffled-label control. Middle and right: the share of this information that the supergroup alone carries, for 35 and 155 supergroups, computed on the classes that share their supergroup with at least one other class.}

\section{Do classes differ in where their sensitivity lies across depth and image space?}\label{depth-space}

The supervised model places most of its sensitivity in blocks 8 to 10, with 17\% of the total in each of blocks 9 and 10. DINO places 35\% of it in block 1 and at most 14\% in any later block. Within this common shape, the class explains 14\% of the variance of an image's mean sensitivity depth in both models, compared with 2\% under shuffled labels. In the supervised model, the mean depth ranges from block 5.4 for carousel, slot and koala to block 7.9 for balloon, parachute and gong.

For the image grid we take each image's 37 most sensitive groups, 0.1\% of all groups, and their per-patch energy. The class explains 10\% of the variance of the vertical centre of this energy in the supervised model and 2.5\% in DINO, against 2\% under shuffled labels. Table lamps and shower caps draw sensitivity from the upper part of the image, unicycles and sarongs from the lower part. The horizontal centre does not depend on the class in either model. The share of energy on the class token depends on the class more strongly than depth does: the class explains 21\% of its variance in the supervised model and 20\% in DINO.

\publicationfigure[\linewidth]
  {figures/kpsa-class-sensitivity-fig3-depth-space.png}
  {/work/kpsa-class-sensitivity/fig3-depth-space.svg}
  {Two heatmaps of sensitivity share per block for 1,000 classes, and a scatter plot of class-mean sensitivity centres on the image grid}
  {Left and middle: the share of sensitivity in each block, averaged over the 50 images of each class, with classes sorted by mean depth. Right: the class-mean centre of the per-patch energy of each image's 37 most sensitive groups on the 14 by 14 patch grid. $\eta^2$ is the share of variance explained by the class.}

\section{Does the class geometry of sensitivity follow the WordNet hierarchy?}\label{hierarchy}

For each block we compare the 1,000 class-mean profiles with one another and correlate their similarity with the depth of the deepest common WordNet ancestor of the two classes. The rank correlation rises from 0.32 in block 1 to 0.49 in block 6 of the supervised model, and from 0.27 to 0.51 in block 7 of DINO. Beyond the middle of the network it falls again, to 0.21 and 0.06 in block 12. The parameters most specific to single classes therefore sit later than those whose use mirrors the full hierarchy. The middle blocks are where the sensitivity of classes is organized coarse-to-fine, consistent with the hierarchical data hypothesis \cite{hdh}.

\publicationfigure[0.62\linewidth]
  {figures/kpsa-class-sensitivity-fig4-hierarchy.png}
  {/work/kpsa-class-sensitivity/fig4-hierarchy.svg}
  {Line plot of the Spearman correlation between class profile similarity and WordNet common-ancestor depth against block, for both models}
  {Spearman correlation, over all pairs of the 1,000 classes, between the similarity of their class-mean profiles in one block and the depth of their deepest common WordNet ancestor. Dotted lines: the same correlation for profiles over all blocks.}

\section{How sparse is the dependence between parameter groups and classes?}\label{sparsity}

We count parameter groups by the exponential of the profile's entropy, the effective number of groups. A single image depends on a median of 2,368 of the 36,864 groups in the supervised model and on 4,823 in DINO, and the top 1\% of groups carry a median 52\% and 35\% of an image's sensitivity. The class mean of 50 images spreads over 5,926 and 8,212 effective groups, so images of one class share part of their groups and add their own.

Counted the other way, a group's sensitivity spreads over the classes according to its depth. Weighted by sensitivity, a group in blocks 1 to 5 serves more than 830 of the 1,000 classes effectively, in both models. In the supervised model this number drops to 27 classes in block 11; in DINO it drops to 47 in block 12. Dependence between parameter groups and classes is therefore dense in early blocks and sparse in late ones, and sparser for the supervised objective.

\publicationfigure[\linewidth]
  {figures/kpsa-class-sensitivity-fig5-sparsity.png}
  {/work/kpsa-class-sensitivity/fig5-sparsity.svg}
  {Histogram of the effective number of parameter groups per image and per class, and a line plot of the effective number of classes per group against block}
  {Left: the effective number of parameter groups, the exponential of the profile entropy, for single images and for class means. Right: the effective number of classes over which a group's sensitivity is spread, as a sensitivity-weighted geometric mean over the groups of a block.}

\section{Implications for KPSA}\label{implications}

A sensitivity atlas retrieves profiles from source inputs that are close to a query. These measurements show where such retrieval has most to gain from a semantic query space: late blocks, where profiles differ between classes and between broad categories, and the block range from 6 to 8, where they follow the class hierarchy. In early blocks, a constant profile already describes most images. Code for these measurements is in \href{https://github.com/Axym-Labs/kpsa}{\texttt{Axym-Labs/kpsa}} under \texttt{src/kpsa/adoption}.

\begin{thebibliography}{99}
\bibitem{kpsa} D. Wiest, A. H. Ariq and A. M. Chowdhury. \href{https://arxiv.org/abs/2610.12737}{Kernelized Parameter Sensitivity Atlases for Amortized Conditional Interpretability}. arXiv:2610.12737, 2026.
\bibitem{steiner2022} A. Steiner et al. \href{https://arxiv.org/abs/2106.10270}{How to train your ViT? Data, Augmentation, and Regularization in Vision Transformers}. \textit{TMLR}, 2022.
\bibitem{caron2021} M. Caron et al. \href{https://arxiv.org/abs/2104.14294}{Emerging Properties in Self-Supervised Vision Transformers}. \textit{ICCV}, 2021.
\bibitem{hdh} D. Wiest. \href{https://axym.org/work/hierarchical-data-hypothesis/}{The Hierarchical Data Hypothesis}. Axym Labs, 2026.
\end{thebibliography}

\end{document}
